\documentclass[12pt,a4paper,UTF8]{ctexart}
\usepackage[a4paper,top=22mm,bottom=22mm,left=22mm,right=22mm,headheight=15pt]{geometry}
\newcommand{\Needspace}[1]{\par\begingroup\dimen0=#1\relax\dimen2=\pagegoal\advance\dimen2 by-\pagetotal\ifdim\dimen0>\dimen2\newpage\fi\endgroup}
\newcommand{\needspace}[1]{\par\begingroup\skip0=\lastskip\vskip-\skip0\dimen0=#1\relax\advance\dimen0 by\skip0\vskip\dimen0\penalty-100\vskip-\dimen0\vskip\skip0\penalty9999\endgroup}
\usepackage{amsmath,amssymb,mathtools,bm}
\usepackage{graphicx,xcolor,booktabs,longtable,array,enumitem,fancyhdr,caption,float,microtype,setspace}
\usepackage[hidelinks,unicode]{hyperref}
\usepackage[most]{tcolorbox}
\definecolor{ContestBlue}{HTML}{17365D}
\hypersetup{pdftitle={第063套},pdfauthor={}}
\setstretch{1.24}
\setlength{\parindent}{2em}
\setlength{\parskip}{0.35em}
\setlength{\emergencystretch}{2em}
\setlist{itemsep=0.35em,topsep=0.4em}
\captionsetup{font=small,labelfont=bf,justification=centering}
\setcounter{secnumdepth}{0}
\pagestyle{fancy}
\fancyhf{}
\fancyhead[L]{\small 第063套}
\fancyfoot[C]{\thepage}
\renewcommand{\headrulewidth}{0.4pt}
\ctexset{
  section={format=\Large\bfseries\color{ContestBlue},beforeskip=1.2em,afterskip=0.7em},
  subsection={format=\large\bfseries\color{ContestBlue},beforeskip=1.0em,afterskip=0.5em},
  subsubsection={format=\normalsize\bfseries,beforeskip=0.8em,afterskip=0.35em}
}
\usepackage{calc}
\newcounter{none}
\let\originalsubsection\subsection
\renewcommand{\subsection}{\needspace{6\baselineskip}\originalsubsection}
\makeatletter
\newcommand{\tightlist}{\setlength{\itemsep}{0pt}\setlength{\parskip}{0pt}}
\newcommand{\maxwidth}{\ifdim\Gin@nat@width>\linewidth\linewidth\else\Gin@nat@width\fi}
\newcommand{\maxheight}{\ifdim\Gin@nat@height>0.72\textheight 0.72\textheight\else\Gin@nat@height\fi}
\newcommand{\pandocbounded}[1]{\begingroup\setkeys{Gin}{width=\maxwidth,height=\maxheight,keepaspectratio}#1\endgroup}
\makeatother
\begin{document}
\begin{titlepage}
  \centering
  \vspace*{58mm}
  {\Huge\bfseries\color{ContestBlue} 第063套\par}
  \vfill
  \vspace*{28mm}
\end{titlepage}
\section{第 1 题}\label{ux7b2c-1-ux9898}

（1）一辆后轮驱动汽车的质量为
\(m=1000\,\mathrm{kg}\)，从静止开始，以恒定切向加速度
\(a=2\,\mathrm{m/s^2}\) 在半径 \(R=500\,\mathrm m\)
的水平圆形车道上行驶。轮胎与路面之间的摩擦因数为
\(\mu=0.5\)。汽车质心位于前、后轮轴的中点，且质心很低，可忽略汽车俯仰和侧倾引起的轮荷转移。求汽车保持在车道上的最长时间。取
\(g=9.8\,\mathrm{m/s^2}\)。

（2）汽车沿圆形轨道做匀速圆周运动。为消除内、外侧车轮压力差，在车上安装一个轮轴与汽车轮轴平行的陀螺轮。陀螺轮角速度与车速成正比，即
\(\omega=kv\)。已知汽车质量为 \(m\)，陀螺轮绕自身轴的转动惯量为
\(I\)，汽车质心离地高度为 \(h\)，求比例系数
\(k\)；同时说明陀螺轮的自转方向应如何选择。

\begin{figure}[H]
\centering
\includegraphics[width=0.53\linewidth,height=\textheight,keepaspectratio,alt={图1-1　汽车在水平圆形车道上的运动几何}]{../figures/question-01.pdf}
\caption*{图1-1　汽车在水平圆形车道上的运动几何}
\end{figure}

\newpage

\section{第 2 题}\label{ux7b2c-2-ux9898}

质量为 \(m\) 的小球位于质量为 \(M=m\)
的山体轨道上。山体左右两侧为直斜面，斜面与竖直方向的夹角均为
\(\alpha=\pi/4\)，每段斜面的竖直高度均为
\(h=1\,\mathrm m\)；底部由一段光滑小圆弧连接，圆弧曲率半径为
\(\rho=10\,\mathrm m\)，圆心角为 \(2\theta=0.02\)。

（1）小球视为质点，山体固定，轨道处处光滑。小球从斜面与圆弧的连接点由静止释放，求往复运动周期。

（2）小球视为质点，山体固定，轨道处处光滑。小球从图示斜面顶端由静止释放，求往复运动周期。

（3）小球视为质点，山体可在光滑水平面上自由平移。小球从图示斜面顶端由静止释放，求系统往复运动周期。

（4）小球改为半径 \(r\ll\rho,h\) 的均质球，绕直径的转动惯量为
\(J=\frac25mr^2\)，运动中始终纯滚动。山体可在光滑水平面上自由平移。均质球从图示斜面顶端由静止释放，求系统往复运动周期。

\begin{figure}[H]
\centering
\includegraphics[width=0.76\linewidth,height=\textheight,keepaspectratio,alt={图2-1　斜面与小圆弧组成的可移动山体轨道}]{../figures/question-02.pdf}
\caption*{图2-1　斜面与小圆弧组成的可移动山体轨道}
\end{figure}

\newpage

\section{第 3 题}\label{ux7b2c-3-ux9898}

如图，两根细杆的质量均为 \(m\)、长度均为 \(l\)。上杆质心到左端 \(A\)
的距离为 \(l_1\)，绕质心且垂直纸面的转动惯量为 \(ma^2\)；下杆质心到左端
\(C\) 的距离为 \(l_2\)，绕质心且垂直纸面的转动惯量为 \(mb^2\)。上杆的
\(A\) 点通过轻质细杆与天花板连接，\(B\)
点通过轻质细绳与天花板连接；上、下两杆之间，\(AC\) 为细绳，\(BD\)
为细杆，各连接点均为铰接。系统初始静止且两杆水平。现剪断 \(B\)
点与天花板之间的细绳，求剪断瞬间 \(BD\) 杆中的轴向作用力。

\begin{figure}[H]
\centering
\includegraphics[width=0.72\linewidth,height=\textheight,keepaspectratio,alt={图3-1　两水平杆与单侧绳---杆约束系统}]{../figures/question-03.pdf}
\caption*{图3-1　两水平杆与单侧绳---杆约束系统}
\end{figure}

\newpage

\section{第 4 题}\label{ux7b2c-4-ux9898}

薄凸透镜 \(L_1\) 与薄凹透镜 \(L_2\) 的焦距分别为 \(f\) 与
\(-2f\)（\(f>0\)），口径半径均为 \(R\)。两透镜平行共轴放置，间距为
\(D=2f\)。

（1）不设光阑，在凸透镜左侧光轴上距 \(L_1\) 为 \(u\)
处放置发光点。光线依次通过两透镜后成实像，求 \(u\) 的最大允许值
\(u_1\)。

（2）在同一条件下，求 \(u\) 的最小允许值 \(u_2\)。

（3）仍在同一条件下，要求最终实像到凹透镜的距离等于发光点到凸透镜的距离，求物距
\(u_3\)。

（4）采用第三问的物距。各透镜均视为理想薄透镜，发光点在三维空间中各向均匀发光。在凸透镜左侧距其
\(f/3\) 处共轴放置圆形光阑，使最终像的亮度为未放光阑时的一半，求光阑半径
\(r\)。

\begin{figure}[H]
\centering
\includegraphics[width=0.78\linewidth,height=\textheight,keepaspectratio,alt={图4-1　凸透镜---凹透镜共轴组合}]{../figures/question-04.pdf}
\caption*{图4-1　凸透镜---凹透镜共轴组合}
\end{figure}

\newpage

\section{第 5 题}\label{ux7b2c-5-ux9898}

如图，一个由电阻和电容构成的无限周期网络以速度 \(u_0\) 沿 \(x\)
轴正方向匀速运动。相邻竖直支路间距、两水平导轨间距均为
\(L\)；每一水平线段的电阻为 \(r\)，每一竖直支路由电阻 \(r\) 与电容 \(C\)
串联组成。空间中存在磁场

\[
\bm B=B_0\cos(\omega t-kx)\,\hat{\bm z}.
\]

\(AC\) 在初始时刻位于 \(x=0\)，随网络运动到
\(x=u_0t\)。原图未画出左侧延伸，且 \(AC\)
画作导线；参考解采用双向无限周期网络的内部支路模型。两种边界解释须分别处理。电容初始电荷趋于零，求暂态衰减后
\(AC\) 支路中的稳态周期电流。

\begin{figure}[H]
\centering
\includegraphics[width=0.867\linewidth,height=\textheight,keepaspectratio,alt={图5-1　平移无限阻容网络与行波磁场}]{../figures/question-05.pdf}
\caption*{图5-1　平移无限阻容网络与行波磁场}
\end{figure}

\newpage

\section{第 6 题}\label{ux7b2c-6-ux9898}

以地心为原点，以从地磁北极指向地磁南极的方向为极轴建立球坐标系。把地球磁场视为磁偶极矩

\[
M=8.1\times10^{22}\,\mathrm{A\,m^2}
\]

产生的偶极子场。

（1）求通过 \(r=R\)、\(\theta=\pi/2\) 点的磁感线方程。

（2）能量 \(E=10\,\mathrm{MeV}\) 的电子在 \(r=R\)、\(\theta=\pi/2\)
附近做回旋运动，求回旋半径 \(a\) 与角频率 \(\omega\) 的数值。已知
\(m_ec^2=0.511\,\mathrm{MeV}\)，\(R=3\times10^7\,\mathrm m\)；按
\(E/(m_ec^2)\approx20\gg1\) 处理。

（3）在第二问运动上叠加沿磁感线方向的小扰动。电子的导引中心将在赤道附近沿磁感线作小振动。求振动角频率
\(\Omega\)，用 \(a\)、\(R\)、\(\omega\)
表示并给出数值。可使用缓变磁场中的第一绝热不变量，即轨道磁矩守恒。

\begin{figure}[H]
\centering
\includegraphics[width=0.536\linewidth,height=\textheight,keepaspectratio,alt={图6-1　地磁偶极场与赤道附近电子回旋}]{../figures/question-06.pdf}
\caption*{图6-1　地磁偶极场与赤道附近电子回旋}
\end{figure}

\newpage

\section{第 7 题}\label{ux7b2c-7-ux9898}

如图，一根质量为 \(M=Nm\)、左端开口而右端封闭的一维钢管中有 \(N\)
个质量均为 \(m\)
的钢球。钢管初始静止在光滑水平地面上，所有碰撞均为一维完全弹性碰撞，钢球初速度均沿管轴方向，不计摩擦。

（1）\(N=4\)，钢球初速度大小均为
\(v_0\)，其中一半向左、一半向右。求末态钢管稳定速度 \(v_{f1}/v_0\)。

（2）保持同样的初始速度分布，令 \(N\to\infty\)，求末态钢管稳定速度
\(v_{f2}/v_0\)。

（3）令 \(N\to\infty\)，把每个钢球换成能量满足 \(h\nu=m_0c^2\)
的光子，其中一半向左、一半向右；光子被钢管封闭端完全弹性反射。求末态钢管稳定速度
\(v_{f3}/c\)。

\begin{figure}[H]
\centering
\includegraphics[width=0.82\linewidth,height=\textheight,keepaspectratio,alt={图7-1　一端开口钢管中的一维弹性碰撞}]{../figures/question-07.pdf}
\caption*{图7-1　一端开口钢管中的一维弹性碰撞}
\end{figure}

\newpage

\section{第 8 题}\label{ux7b2c-8-ux9898}

两条光滑、足够长的水平导轨间距为 \(l\)，空间中存在竖直向下的匀强磁场
\(B\)。一根质量为 \(m\)、等效电感为 \(L\)
的金属杆垂直搭在导轨上，初始静止。电路中串联一个初始电压为
\(\lambda U_0\) 的电容器 \(C\)，并接有两个反向并联、阈值电压均为 \(U_0\)
的理想二极管。闭合开关并等待足够长时间，求金属杆的末态速度。

\begin{figure}[H]
\centering
\includegraphics[width=0.86\linewidth,height=\textheight,keepaspectratio,alt={图8-1　含电感杆、电容与双阈值二极管的磁耦合电路}]{../figures/question-08.pdf}
\caption*{图8-1　含电感杆、电容与双阈值二极管的磁耦合电路}
\end{figure}
\end{document}
